\documentclass[12pt]{article}
\usepackage[utf8]{inputenc}
\usepackage[T1]{fontenc}
\usepackage{hyperref}
\title{Avava Language Resources for K-12 Curriculum Development}
\author{Compiled by Claude Sonnet 5 (claude-sonnet-5) for Andr\'as Kornai}
\date{}
\begin{document}
\maketitle

\section{Introduction}
Avava (ISO 639-3 \texttt{tmb}, Glottocode katb1237; also called Katbol, Tembimbe-Katbol, Bangsa') is an Oceanic language of central Malekula, Vanuatu, with about 700 speakers in villages such as Tisvel, Khatbol, Taremp and Tembimbi. Timbembe and Nevaar dialects are still spoken. Children reportedly learn Avava alongside Bislama. Vitality assessments differ between sources.

\section{Primary Educational Resources}
No primers, textbooks or school programmes were found.

\section{Linguistic and Reference Materials}
\begin{itemize}
\item Terry Crowley (2006), \emph{The Avava Language of Central Malakula (Vanuatu)}, Pacific Linguistics, the main grammar (Glottolog reference): \url{https://glottolog.org/resource/reference/id/151851}
\item ANU Open Research repository file surfaced in the search (contents not verified; possibly the Crowley grammar): \url{https://openresearch-repository.anu.edu.au/bitstreams/80607528-418e-4d9a-af23-5088898bb720/download}
\item Wikipedia article on Avava: \url{https://en.wikipedia.org/wiki/Avava_language}
\item Austronesian Basic Vocabulary Database entry: \url{https://abvd.eva.mpg.de/austronesian/language.php?id=456}
\item Omniglot page on Avava: \url{https://www.omniglot.com/writing/avava.htm}
\item Endangered Languages Project entry: \url{https://endangeredlanguages.com/lang/10209}
\end{itemize}

\section{Religious and Community Resources}
None found.

\section{Digital and Audio Resources}
None found beyond the wordlist database above.

\section{Curriculum Development Context}
Documentation is limited to Crowley's grammar and short wordlists. A curriculum would need a community-developed orthography and locally produced materials, building on the grammar and with Bislama and English as bridging languages.
\end{document}